\chapter{Preface}

This series has a simple ambition: a single, coherent course in quantitative
finance, from the plumbing of the markets to the running of the firm, written
for readers anywhere in the world. It covers every major market, the fast and
the slow ends of trading, the bank and the asset manager, the mathematics and
the machines.

The reader is assumed to program fluently in C++, Rust and Python and to have
the mathematical culture of a master's degree. Nothing here re-teaches a
language or undergraduate mathematics; everything here is the layer on top.

The style is concise and rigorous: lessons built from definitions, examples,
propositions, methods and figures, with careful reasoning wherever space
allows. Results stated without full derivation are explicitly marked as
admitted. Each chapter then makes the reader work: a guided tutorial with
code, a build that adds a piece to a miniature trading firm assembled across
the series, graded exercises, one long weekend problem, and a set of
interview questions. Full solutions are collected at the end of the book.

\section*{How to read this book}

Statements are numbered within each chapter and color-coded: definitions in
blue, theorems in red, propositions and their relatives in orange, and
methods --- short recipes for standard tasks --- in green. Strategy files are
purple, tutorials and builds slate, interview questions teal. Ochre boxes
headed ``As of'' hold facts that change: they carry the month in which they
were checked against a public document. Exercises and interview questions are
graded from ($\star$) for direct applications to ($\star\star\star$) for
harder problems.

Every listing is an excerpt of a tested file in the book's repository; its
path is printed under it. Type the tutorials, do the builds: the series is
written to be worked through, not read.

\section*{Contributing}

The source of this book lives in a public git repository:
\begin{center}
\url{https://github.com/bvirrion/one-quant-book}
\end{center}
Corrections, better sources and additional exercises are
welcome; see \texttt{CONTRIBUTING.md} at the root of the repository.
